% Generated draft: requires mathematical and source review.
\documentclass[11pt]{article}
\usepackage[margin=25mm]{geometry}
\usepackage{fontspec}
\setmainfont{Latin Modern Roman}
\usepackage{amsmath,amssymb,mathtools,unicode-math}
\setmathfont{Latin Modern Math}
\usepackage{enumitem,needspace,xurl,hyperref}
\setlength{\parindent}{0pt}
\setlength{\parskip}{5pt}
\setlength{\emergencystretch}{3em}
\newcommand{\N}{\mathbb{N}}
\newcommand{\Z}{\mathbb{Z}}
\newcommand{\Q}{\mathbb{Q}}
\newcommand{\R}{\mathbb{R}}
\newcommand{\C}{\mathbb{C}}
\newcommand{\F}{\mathbb{F}}
\newcommand{\A}{\mathbb{A}}
\newcommand{\PP}{\mathbb{P}}
\newcommand{\E}{\mathbb{E}}
\newcommand{\eps}{\varepsilon}
\newcommand{\dd}{\,\mathrm{d}}
\begin{document}
\section*{UnsolvedMath ID: 1100302}
\textbf{Source problem code:} catalogue entry
\subsection{Definitions and mathematical statement}
Let $\mathbb H^n$ have its standard metric of curvature $-1$. For a Kleinian group $H\leq\operatorname{Isom}(\mathbb H^n)$, let $\Lambda H\subseteq\partial\mathbb H^n$ be its limit set and, when the convex-hull construction is defined, write
\[
\operatorname{Core}(H)=\operatorname{Hull}(\Lambda H)/H,
\qquad v_n(H)=\operatorname{Vol}_n(\operatorname{Core}(H)).
\]
Here the hull is taken inside $\mathbb H^n$, and the volume is the ambient $n$-dimensional hyperbolic volume, not an intrinsic lower-dimensional volume. Elementary groups and zero-volume cores require compatible conventions if included.

In dimension $2$, the motivating formula for a nontrivial finitely generated free Fuchsian group $A$ in the source's setting is
\[
v_2(A)=2\pi\bigl(\operatorname{rk}(A)-1\bigr).
\]
Consequently its stated Hanna Neumann inequality is
\[
2\pi\,v_2(A\cap B)\leq v_2(A)v_2(B)
\]
for nontrivial free subgroups $A,B$ of a Fuchsian group, with zero core area for trivial intersections understood if necessary. The source's opening definition says ``area$(G)$'' but then uses $\mathbb H^2/A$; the notation $v_2(A)$ above consistently uses the group whose core is measured.

Swarup asks for analogous inequalities for torsion-free, not necessarily free, quasiconvex subgroups $A,B$ of a quasiconvex Kleinian group $\Gamma\leq\operatorname{Isom}(\mathbb H^n)$, replacing area by $v_n$. A proposed inequality has the form
\[
c\,v_n(A\cap B)\leq v_n(A)v_n(B),\qquad c>0,
\]
possibly with a constant worse than $2\pi$. The source does not specify the permitted dependence of $c$. A universal dimension-dependent version would require $\forall n\geq3\ \exists c_n>0\ \forall\Gamma,A,B$ as above; a version allowing $c=c(\Gamma)$ is different. Neither quantifier choice is silently imposed here. Nor does the record define whether ``quasiconvex Kleinian group'' means a group with a quasiconvex orbit in $\mathbb H^n$ (the usual convex-cocompact setting), or another convention; admissibility and finiteness of the core volumes must be fixed accordingly.

The supplied research summary reports the free Fuchsian case as the Friedman--Mineyev theorem. It further claims a sharp-$2\pi$ extension to arbitrary torsion-free finitely generated Fuchsian subgroups by index/Euler-characteristic arguments in an omitted ``Work done'' section; that argument is not supplied here. It explicitly leaves the $n\geq3$ problem unresolved and says that its attempted $\ell^2$-Betti-number and rank-based reductions fail: the first available submultiplicativity result concerns graphs and does not measure core volume, and core volume is topologically unbounded in dimension $3$. These are limitations reported by the summary, not new theorems established in this draft.
\subsection{Short English statement}
The question seeks a higher-dimensional convex-core-volume analogue of the Hanna Neumann intersection inequality, allowing a weaker positive constant and nonfree torsion-free quasiconvex subgroups. The constant's allowed dependence and the precise Kleinian-group conventions are unspecified. The supplied summary leaves dimensions at least $3$ open and gives no accessible proof for its additional claimed Fuchsian extension.
\subsection{Sources}
\begin{itemize}
\item Local UnsolvedMath registry, permanent ID 1100302, exact source code catalogue entry, ``Questions in Geometric Group Theory --- Q 3.2''; status partially solved, classification PARTIAL-PROGRESS.
\item Bestvina, \emph{Questions in Geometric Group Theory} (2004), Question 3.2, PDF page 10, attributed to Swarup, \url{https://www.math.utah.edu/~bestvina/eprints/questions-updated.pdf}. The record identifies manual transcription of a public author-hosted PDF, accessed July 29, 2026, and says the rendered source presents it as open with current status not independently verified.
\item The supplied research summary names the Friedman--Mineyev theorem for free subgroups, without providing bibliographic details. Its asserted extension to all torsion-free finitely generated Fuchsian subgroups refers to ``Work done,'' item 1, which is absent from this record; no authoritative citation or proof of that extension is supplied.
\item The research summary explicitly reports no solution or counterexample for $n\geq3$. Its claims about failed approaches, intersection of convex hulls of limit sets, and a thin/thick dichotomy refer to preceding material not supplied here. Its statement that no literature addresses the higher-dimensional question is a supplied assessment, not a verified exhaustive literature search.
\end{itemize}
\end{document}
